\documentclass{article}
\usepackage[utf8]{inputenc}
\usepackage{amsmath, amssymb}
\usepackage{siunitx}
\usepackage{booktabs}
\title{Detaljerte beregninger for EMP-generator (revidert)}
\author{}
\date{}

\begin{document}
\maketitle

\section{Innledning}
Dette dokumentet inneholder teoretiske beregninger for designet. Alle verdier er idealiserte og må verifiseres eksperimentelt. Vi skiller mellom matematisk korrekte utledninger, antatte parametere og estimater med usikkerhet.

\section{Hovedkondensator og spole}

\subsection{Spesifikasjoner}
\begin{itemize}
\item Kondensator: $C = 1000\,\mu\text{F}$, merkespenning $450\,\text{V}$ (lades til $V_0 = 400\,\text{V}$ for sikkerhetsmargin).
\item Spole: Flat spiral, $N = 6$ vindinger, indre diameter $20\,\text{mm}$, ytre diameter $100\,\text{mm}$, tråd: $10\,\text{AWG}$ ($2.588\,\text{mm}$ diameter).
\end{itemize}

\subsection{Induktans for flat spiralspole}
Wheelers formel for en flat spiralspole i luft, med dimensjoner i tommer og resultat i µH, er:
\begin{equation}
L = \frac{r^2 N^2}{8r + 11w},
\end{equation}
der $r$ er middelradius i tommer, $w$ er viklingens bredde i tommer.

Konvertering: $r = 30\,\text{mm} = 1.181\,\text{tommer}$, $w = 40\,\text{mm} = 1.575\,\text{tommer}$.
\begin{equation}
L = \frac{(1.181)^2 \cdot 6^2}{8 \cdot 1.181 + 11 \cdot 1.575} = \frac{1.395 \cdot 36}{9.448 + 17.325} = \frac{50.22}{26.773} \approx 1.876\,\mu\text{H}.
\end{equation}
Vi bruker $L = 1.9\,\mu\text{H}$ som et teoretisk estimat. Målt verdi kan avvike.

\subsection{Energi lagret i kondensator}
\begin{equation}
E_0 = \frac{1}{2} C V_0^2 = 0.5 \cdot 1000 \times 10^{-6} \cdot (400)^2 = 80\,\text{J}.
\end{equation}

\section{Toppstrøm og pulsform}

\subsection{Karakteristisk impedans}
\begin{equation}
Z_0 = \sqrt{\frac{L}{C}} = \sqrt{\frac{1.9 \times 10^{-6}}{1000 \times 10^{-6}}} = \sqrt{1.9 \times 10^{-3}} \approx 0.0436\,\Omega.
\end{equation}

\subsection{Toppstrøm (ideell, tapsfri)}
\begin{equation}
I_{\text{peak, ideal}} = \frac{V_0}{Z_0} = \frac{400}{0.0436} \approx 9\,174\,\text{A}.
\end{equation}

\subsection{Demping og reell toppstrøm}
Total motstand $R$ i kretsen er dominert av ESR i kondensator, ledningsmotstand og lysbuemotstand. Vi antar $R \approx 80\,\text{m}\Omega$ basert på typiske verdier, men dette er ikke målt. Dempingsfaktoren:
\begin{equation}
\alpha = \frac{R}{2L} = \frac{0.08}{2 \cdot 1.9 \times 10^{-6}} \approx 21\,050\,\text{s}^{-1}.
\end{equation}
Resonansvinkelfrekvens:
\begin{equation}
\omega_0 = \frac{1}{\sqrt{LC}} = \frac{1}{\sqrt{1.9 \times 10^{-6} \cdot 1000 \times 10^{-6}}} = \frac{1}{\sqrt{1.9 \times 10^{-9}}} \approx 22\,940\,\text{rad/s}.
\end{equation}
Frekvens: $f = \omega_0/(2\pi) \approx 3.65\,\text{kHz}$.
Siden $\alpha < \omega_0$, er kretsen underdempet. Toppstrømmen reduseres noe av dempingen; et realistisk estimat er $8\,000$--$9\,000$ A, men dette er ikke verifisert.

\subsection{Pulsvarighet}
Periode: $T = 1/f \approx 274\,\mu\text{s}$. Tidskonstant for innhyllingskurven: $\tau = 1/\alpha \approx 47.5\,\mu\text{s}$. Etter $5\tau \approx 240\,\mu\text{s}$ er svingningene nesten borte.

\section{Magnetfelt}

\subsection{Felt i sentrum av spolen}
\begin{equation}
B = \frac{\mu_0 N I}{2 r}.
\end{equation}
Med $N=6$, $r=0.05\,\text{m}$, $I=9\,000\,\text{A}$:
\begin{equation}
B = \frac{4\pi \times 10^{-7} \cdot 6 \cdot 9000}{2 \cdot 0.05} \approx 0.679\,\text{T}.
\end{equation}

\subsection{Feltets avtagning med avstand}
For $d = 10\,\text{cm}$:
\begin{equation}
B(0.1) \approx 0.679 \left( \frac{0.05}{\sqrt{0.05^2 + 0.1^2}} \right)^3 \approx 0.679 \cdot 0.089 \approx 0.060\,\text{T}.
\end{equation}

\section{Indusert spenning i test-sløyfe}
Med $N_p=5$, $A_p=3.14\times10^{-4}\,\text{m}^2$, $\Delta t \approx T/4 = 68.5\,\mu\text{s}$:
\begin{equation}
\mathcal{E} \approx -5 \cdot 3.14\times10^{-4} \cdot \frac{0.060}{68.5\times10^{-6}} \approx -1.38\,\text{V}.
\end{equation}
Dette er et teoretisk estimat. Faktisk $dB/dt$ kan være høyere.

\section{Energibalanse}

\subsection{Energibevaring under utladning}
Den totale energien $E_0 = 80\,\text{J}$ er opprinnelig lagret i kondensatorens elektriske felt. Under utladningen omdannes denne energien kontinuerlig. Til enhver tid gjelder:

\begin{equation}
E_0 = E_C(t) + E_L(t) + E_{\text{diss}}(t) + E_{\text{utstrålt}}(t),
\end{equation}

der:
\begin{itemize}
\item $E_C(t) = \frac{1}{2} C v_C(t)^2$ er gjenværende elektrisk feltenergi i kondensatoren,
\item $E_L(t) = \frac{1}{2} L i(t)^2$ er magnetisk feltenergi midlertidig lagret i spolen,
\item $E_{\text{diss}}(t) = \int_0^t i^2(\tau) R\,d\tau$ er energi omsatt til varme i ohmsk motstand (ledninger, spole, ESR),
\item $E_{\text{utstrålt}}(t)$ er energi forlatt systemet som elektromagnetisk stråling (liten ved disse frekvensene).
\end{itemize}

I tillegg kommer energitap i lysbuen, som ikke er en enkel ohmsk prosess. Denne modellen neglisjerer lysbuens ikke-linearitet og strålingstap.

Når utladningen er fullført ($t \to \infty$), er $E_C = 0$, $E_L = 0$, og all energi er omdannet til varme og stråling. Det er **ikke korrekt** å tilordne hele $E_0$ til oppvarming av én enkelt komponent; energien fordeles mellom flere tapsmekanismer.

\subsection{Termiske betraktninger (hypotetiske)}
Følgende beregninger er **kun illustrative eksempler** og forutsetter vilkårlige fordelingsandeler. De kan ikke brukes til å fastslå faktisk temperaturøkning eller komponentsikkerhet.

\subsubsection{Eksempel: Spole}
Dersom 20\,\% av $E_0$ ($16\,\text{J}$) går til oppvarming av spolen (masse $\approx 89\,\text{g}$, varmekapasitet $385\,\text{J/(kg·K)}$):
\begin{equation}
\Delta T \approx \frac{16}{0.089 \cdot 385} \approx 0.47\,\text{K}.
\end{equation}

\subsubsection{Eksempel: Kondensator}
Dersom 30\,\% av $E_0$ ($24\,\text{J}$) tapes i kondensatorens ESR, og kondensatoren har masse $\approx 50\,\text{g}$ og varmekapasitet $\approx 50\,\text{J/K}$:
\begin{equation}
\Delta T \approx \frac{24}{50} = 0.48\,\text{K}.
\end{equation}

\subsection{Konklusjon termisk}
De faktiske temperaturøkningene avhenger av den ukjente energifordelingen. Lokale varmepunkter, gjentatte pulser og komponentenes transienttoleranse er ikke vurdert. **Sikker drift kan ikke garanteres uten eksperimentell verifikasjon.**

\section{Komponentpåkjenninger}

\subsection{MOSFET i ladekrets}
Under avslag kan drain-spenningen nå opp til $2 \times V_{\text{bat}} + V_{\text{reflektert}}$. Med $V_{\text{bat}} = 54\,\text{V}$ og et viklingsforhold $n = 200/10 = 20$, er reflektert spenning $V_{\text{ut}}/n = 400/20 = 20\,\text{V}$. Total spenning: $2 \cdot 54 + 20 = 128\,\text{V}$. IRF840 tåler $500\,\text{V}$, så marginen er god.

\subsection{Dioder}
UF4007 tåler $1000\,\text{V}$ repeterende. I verste fall kan utgangsspenningen nå $450\,\text{V}$, så marginen er over 2x.

\subsection{Gnistgap}
Wolframelektroder har et smeltepunkt på $3422\,^{\circ}\text{C}$ og tåler ekstremt høye pulsstrømmer uten å erodere nevneverdig.

\section{Konklusjon}
Alle beregninger er teoretiske. Fysisk prototype må bygges og måles for å bekrefte ytelse og sikkerhet.
\end{document}
